% Responsable: Anthony · Ronda 1 (viernes 13, hasta 21.00)
% Para cada concepto: (i) definición formal con cita IEEE comprimida, (ii) explicación propia de mínimo 3 oraciones, (iii) ejemplo del Sistema de gestión para un restaurante con IA.
\subsection{Tipos de sistemas}
La clasificación de un sistema permite identificar su propósito, su entorno
de operación y la relación que mantiene con las personas y otros elementos
externos. En Ingeniería de Requisitos esta caracterización ayuda a establecer
el alcance del producto y a comprender las necesidades de los stakeholders
antes de especificar las funciones concretas del software
\cite{ieee29148,pohl2025}.

\subsubsection{Sistema de información de gestión}

\textbf{(i) Definición formal.}
Un sistema de información permite registrar, almacenar, procesar y comunicar
datos con el propósito de apoyar las actividades de una organización. Desde
la perspectiva de Ingeniería de Requisitos, sus funciones deben responder a
las necesidades de los stakeholders y al contexto en el que será utilizado
\cite{ieee29148,pohl2025}.

\textbf{(ii) Explicación propia.}
Este tipo de sistema trabaja principalmente con información generada durante
las actividades de una organización. Su objetivo no consiste únicamente en
guardar datos, sino también en procesarlos para facilitar las operaciones y
apoyar la toma de decisiones. Por esta razón, sus requisitos deben establecer
qué información se registra, quién puede utilizarla y qué resultados debe
proporcionar el sistema.

\textbf{(iii) Aplicación al PFC.}
El Sistema de gestión para un restaurante con IA desarrollado para
El Balcón Restaurant, en Quevedo, se clasifica principalmente como un sistema
de información de gestión. La aplicación registrará y procesará información
relacionada con pedidos, reservas, menú, inventario, ventas y cierre de caja.
Estos datos permitirán apoyar las actividades diarias del restaurante y
servirán también como fuente para las funcionalidades de inteligencia
artificial definidas en el proyecto.

\subsubsection{Sistema sociotécnico}

\textbf{(i) Definición formal.}
El análisis de requisitos debe considerar no solamente los componentes de
software, sino también a las personas, los procesos organizacionales y el
entorno en el que funcionará el sistema. Esta interacción entre elementos
humanos y tecnológicos caracteriza a los sistemas sociotécnicos
\cite{pohl2025,swebok2024}.

\textbf{(ii) Explicación propia.}
El funcionamiento correcto de una aplicación empresarial no depende
exclusivamente del código desarrollado. Los usuarios realizan actividades,
siguen procedimientos y toman decisiones que influyen directamente en los
resultados obtenidos mediante el sistema. Por ello, durante la Ingeniería de
Requisitos es necesario comprender tanto la parte tecnológica como la forma
real en que trabajan las personas dentro de la organización.

\textbf{(iii) Aplicación al PFC.}
El sistema de El Balcón Restaurant posee una naturaleza sociotécnica porque
en sus procesos intervienen el propietario, el jefe de cocina, los meseros,
el cajero y los clientes. Por ejemplo, el mesero registra un pedido, cocina
recibe la información para preparar los platos y posteriormente el cajero
utiliza los datos registrados para realizar el cobro. El resultado final
depende de la interacción entre los usuarios, los procedimientos internos
del restaurante y la aplicación.

\subsubsection{Sistema con componente de inteligencia artificial}

\textbf{(i) Definición formal.}
Dentro del proyecto, la inteligencia artificial constituye un componente
integrado al sistema de información. Las capacidades incorporadas al software
deben originarse en necesidades identificadas y expresarse mediante requisitos
que puedan ser documentados y posteriormente validados
\cite{ieee29148,pohl2025}.

\textbf{(ii) Explicación propia.}
La incorporación de inteligencia artificial no cambia el propósito principal
del sistema, que continúa siendo gestionar información y procesos del
restaurante. La IA utiliza los datos disponibles para generar resultados que
sirven como apoyo a determinadas decisiones. En este proyecto sus resultados
se consideran recomendaciones, por lo que las personas mantienen la
responsabilidad sobre las decisiones finales.

\textbf{(iii) Aplicación al PFC.}
En El Balcón Restaurant se han definido dos funciones principales de
inteligencia artificial. La primera permitirá estimar la demanda diaria de
cada plato utilizando información histórica de ventas, con el objetivo de
apoyar la planificación de compras de insumos. La segunda permitirá recomendar
platos a los clientes tomando como referencia sus pedidos anteriores. En ambos
casos la inteligencia artificial funciona como apoyo y no como sustituto de
la decisión humana.

\subsubsection{Clasificación del PFC}

De acuerdo con las características identificadas, el proyecto se considera
principalmente un \textbf{sistema de información de gestión con un componente
de inteligencia artificial}. También presenta una naturaleza sociotécnica,
debido a que su funcionamiento depende de la interacción entre el software,
los usuarios y los procedimientos internos de El Balcón Restaurant.

El sistema no se considera embebido, ya que su función principal no consiste
en controlar directamente dispositivos físicos. Tampoco se clasifica como un
sistema crítico de seguridad, puesto que las funciones de inteligencia
artificial generan estimaciones y recomendaciones que posteriormente son
evaluadas por las personas responsables. Esta clasificación permite orientar
los requisitos hacia la gestión de información, la operación del restaurante
y el apoyo a la toma de decisiones.
